(1) 相互作用绘景一阶振幅 $c_E^{(1)}(t)=-(i/\hbar)V(E)\int_0^t e^{i(E-E_0)s/\hbar}ds$，积分后可写 $c_E^{(1)}(t)=V(E)[1-e^{i(E-E_0)t/\hbar}]/(E-E_0)$。在能量相等点应取可去奇点极限 $-iVt/\hbar$。

(2) 能量窗口内概率密度为 $|c_E|^2=4|V(E)|^2\sin^2[(E-E_0)t/(2\hbar)]/(E-E_0)^2$。若连续谱态密度 $\rho(E)$ 缓变，则分布核在长时间满足 $\frac1t4\sin^2(\Delta Et/(2\hbar))/\Delta E^2\to(2\pi/\hbar)\delta(\Delta E)$，从而得到 $\boxed{\Gamma=2\pi|V(E_0)|^2\rho(E_0)/\hbar}$，单位为逆时间。

(3) 在弱耦合且无记忆近似中，离散态振幅的概率满足 $\dot P_d=-\Gamma P_d$，因此 $\boxed{P_d(t)=e^{-\Gamma t},\quad\tau=1/\Gamma}$。相应指数振幅Fourier变换的强度线形近似为 $[ (E-E_0)^2+(\hbar\Gamma/2)^2]^{-1}$，能量全宽半高是 $\boxed{\Delta E_{\rm FWHM}=\hbar\Gamma}$。

(4) 严格量子短时演化对任意具有有限能量方差的初态有 $P_d(t)=1-(\Delta H)^2t^2/\hbar^2+O(t^3)$，故初始瞬时衰减斜率为零，不可能从 $t=0$ 完全符合指数式。强耦合、有限谱边、带隙、强能量依赖态密度和长时间幂律尾部也会破坏简单Markov黄金律。
